% ID: SEL-D0010-BB-E04
% ESTADO: cotejada. Segunda pasada visual completa y borrador compilado revisados.
% Fuente: Selectividad/matematicas_aplicadas/EXTRA_TITULAR_MAT-APLIC-CCSS-JULIO-IMPRESO.pdf
% SHA-256: 548f60ee334516c7fa918dba3c28fb3ea5fe753cf76eb05402d7b8acb5a3b560
% Página: 1; bloque B; ejercicio original 4.

El número de diagnosticados de COVID-19 por PCR en Andalucía, medido
en miles de personas, se aproxima por la siguiente función:
\[
f(t)=\begin{cases}
-t^2+2t-0.3 & \text{si }0.2\leq t\leq 1.8\\
0.1t-0.12 & \text{si }1.8<t\leq 5\\
-0.5t^2+8.3t-28.62 & \text{si }5<t\leq 10
\end{cases}
\]
donde $t$ es el tiempo, medido en meses, a partir del inicio de conteo en
el mes de marzo de 2020.
\begin{enumerate}[label=\alph*)]
  \item \textbf{(1.5 puntos)} Estudie la continuidad y derivabilidad de la
  función $f$ en su dominio.
  \item \textbf{(1 punto)} \textquestiondown{}En qué instante o instantes
  es máximo el número de diagnosticados? \textquestiondown{}Cuál es ese número?
\end{enumerate}
